% Cover letter for the JCIM submission. Build: latexmk -pdf cover_letter.tex
\documentclass[11pt,letterpaper]{article}
\usepackage[T1]{fontenc}
\usepackage[utf8]{inputenc}
\usepackage[margin=0.9in,top=0.7in,bottom=0.7in]{geometry}
\usepackage[colorlinks=true,allcolors=blue]{hyperref}
\usepackage{parskip}
\setlength{\parskip}{0.55em}
\pagestyle{empty}

\begin{document}

\noindent
Isman Kurniawan \\
School of Computing, Telkom University \\
Bandung, Indonesia \\
\texttt{isman.krn@gmail.com} \\[0.4em]
\today

\medskip
\noindent
To the Editors, \emph{Journal of Chemical Information and Modeling}

\medskip
\noindent
Dear Editors,

We submit for your consideration an Article entitled \textbf{``Activity cliffs on human
acetylcholinesterase: a pre-specified descriptor escalation''}.

It is widely held that two-dimensional fingerprints fail on activity cliffs because
they cannot see the three-dimensional detail separating the two members of a pair, and
that supplying that detail should help. We tested that directly, escalating four
structure-based representations against the same fingerprint baseline on the same
frozen splits, and measured in advance how large a difference the design could have
resolved. None was resolvably better on any of three framings of the endpoint, and on
per-ligand potency the molecular dynamics descriptors were resolvably worse. The result
is bounded rather than open, which is what separates it from an underpowered null.

Three findings carry beyond this target. A replicate test separates persistence from
amplitude sharply, so a trajectory of affordable length reproduces how long a contact
is held and not how far things move. A descriptor filter refitted at each value of a
varied parameter manufactures a trend that a single fixed selection removes, which
exposes any study that re-selects features as it varies a parameter. And the
descriptors that survive reproducibility are available from 15 ns of simulation rather
than the 50 the original design called for.

We would like to address three points of the Journal's guidance on QSAR submissions
before a reader reaches them. First, this is not a docking study against a single
target: the docking serves a comparison between representations, and the quantity under
test is what physical detail adds to a matched ligand-based baseline. Second, the
high ratio of descriptors to data points sits in the fingerprint baseline, scored as a
practitioner would score it, and we report a matched-dimensionality arm alongside it
precisely so that feature count can be separated from representation. Third, we are not
proposing new descriptors and claim no superiority for them; what is offered is the
comparison and the detection limit that makes its outcome readable, and the descriptor
machinery is standard and attributed as such.

Acetylcholinesterase was chosen because the question needs a target supplying activity
cliffs in quantity and a crystallographic human receptor, and because its gorge gates
conformationally on evidence predating this work. A comparison that finds no advantage
where one is most expected says more than one that finds none where none was expected.

All molecular structures appear in the Supporting Information, and all analysis code
and derived data are in a public repository archived with a DOI, sufficient to
reproduce every number without rerunning a simulation.

The manuscript is original, is not under consideration elsewhere, and has been approved
by all authors. The authors declare no competing financial interest, and the work
received no external funding.

\medskip
\noindent
Yours sincerely,

\medskip
\noindent
Isman Kurniawan, on behalf of all authors

\end{document}
